\section{Formalisasi Kalkulus Predikat: Term, Proposisi, Kalimat, dan Ekspresi}

Pada section sebelumnya, kita telah mempelajari predikat, variabel, dan kuantifier secara intuitif. Section ini akan memformalkan konsep-konsep tersebut dengan mendefinisikan secara ketat unsur-unsur pembentuk bahasa kalkulus predikat.

\subsection{Simbol-Simbol dalam Kalkulus Predikat}

Kalimat dalam kalkulus predikat dibangun dari simbol-simbol berikut:

\begin{konsep}
\begin{enumerate}
  \item \textbf{Simbol Kebenaran}: \code{true} dan \code{false}
  \item \textbf{Simbol Konstanta}: $a, b, c, a_1, b_1, \ldots$ — mewakili objek tertentu dalam domain
  \item \textbf{Simbol Variabel}: $x, y, z, x_1, x_2, \ldots$ — mewakili objek yang belum ditentukan
  \item \textbf{Simbol Fungsi}: $f, g, h, g_1, f_1, h_1, \ldots$ — setiap simbol fungsi mempunyai \emph{arity} yang menyatakan banyaknya parameter/argumen yang harus dipenuhi
  \item \textbf{Simbol Predikat} (menyatakan relasi): $p, q, r, s, p_1, q_1, \ldots$ — setiap simbol predikat juga memiliki arity
\end{enumerate}
\end{konsep}

\subsection{Definisi Term}

\begin{konsep}
\logic{Term} adalah sebuah ekspresi yang menyatakan objek. Term dibangun berdasarkan aturan-aturan berikut:
\begin{enumerate}
  \item Semua \textbf{konstanta} adalah term
  \item Semua \textbf{variabel} adalah term
  \item Jika $t_1, t_2, \ldots, t_n$ adalah term ($n \geq 1$) dan $f$ adalah fungsi dengan arity $= n$, maka $f(t_1, t_2, \ldots, t_n)$ adalah term
  \item Jika $A$ adalah kalimat, sedangkan $s$ dan $t$ adalah term, maka kondisional ``\code{if} $A$ \code{then} $s$ \code{else} $t$'' adalah term
\end{enumerate}
\end{konsep}

\begin{contoh}
\begin{enumerate}
  \item $f(a, x)$ adalah term, karena:
  \begin{itemize}
    \item $a$ adalah simbol konstanta, dan semua konstanta adalah term
    \item $x$ adalah simbol variabel, dan semua variabel adalah term
    \item $f$ adalah simbol fungsi dan $f(a,x)$ memenuhi aturan ke-3
  \end{itemize}

  \item $g(x, f(a, x))$ adalah term, karena:
  \begin{itemize}
    \item $a$ adalah konstanta (term), $x$ adalah variabel (term)
    \item $f(a, x)$ adalah term (dari contoh sebelumnya)
    \item $g$ adalah simbol fungsi dengan arity 2, dan $x$ serta $f(a,x)$ keduanya term
  \end{itemize}
\end{enumerate}
\end{contoh}

\subsection{Definisi Proposisi (Atom)}

\begin{konsep}
\logic{Proposisi} (atom) digunakan untuk merepresentasikan relasi antar objek. Proposisi dibangun berdasarkan aturan berikut:
\begin{enumerate}
  \item Simbol kebenaran (\code{true}, \code{false}) adalah proposisi
  \item Jika $t_1, t_2, \ldots, t_n$ adalah term dan $p$ adalah simbol predikat dengan arity $n$, maka $p(t_1, t_2, \ldots, t_n)$ adalah proposisi
\end{enumerate}
\end{konsep}

\begin{contoh}
$p(a, x, f(a, x))$ adalah proposisi, karena:
\begin{itemize}
  \item $a$ adalah simbol konstanta (term)
  \item $x$ adalah simbol variabel (term)
  \item $f(a, x)$ adalah term (fungsi diterapkan pada term)
  \item $p$ adalah simbol predikat dengan arity 3 (3-ary)
\end{itemize}
\end{contoh}

\subsection{Definisi Kalimat}

\begin{konsep}
\logic{Kalimat} dalam kalkulus predikat dibangun dengan aturan:
\begin{enumerate}
  \item Setiap proposisi adalah kalimat
  \item Jika $A$, $B$, $C$ adalah kalimat, maka:
  \begin{enumerate}
    \item \textbf{Negasi}: (\code{not} $A$) adalah kalimat
    \item \textbf{Konjungsi}: ($A$ \code{and} $B$) adalah kalimat
    \item \textbf{Disjungsi}: ($A$ \code{or} $B$) adalah kalimat
    \item \textbf{Implikasi}: (\code{if} $A$ \code{then} $B$) adalah kalimat
    \item \textbf{Ekuivalensi}: ($A$ \code{iff} $B$) adalah kalimat
    \item \textbf{Kondisional}: (\code{if} $A$ \code{then} $B$ \code{else} $C$) adalah kalimat
  \end{enumerate}
  \item Jika $A$ adalah kalimat dan $x$ adalah variabel, maka:
  \begin{enumerate}
    \item $(\forall x)\, A$ adalah kalimat
    \item $(\exists x)\, A$ adalah kalimat
  \end{enumerate}
\end{enumerate}
\end{konsep}

\begin{contoh}
\begin{enumerate}
  \item \code{if} $(\forall x)\, p(a, b, x)$ \code{then} $g(y)$ \code{else} $f(a, y)$ adalah \textbf{term}, karena:
  \begin{itemize}
    \item $(\forall x)\, p(a, b, x)$ adalah kalimat
    \item $g(y)$ dan $f(a, y)$ adalah term
    \item Maka kondisional \code{if} kalimat \code{then} term \code{else} term adalah term
  \end{itemize}

  \item \code{if} $(\forall x)\, p(a, b, x)$ \code{then} $(\exists y)\, q(y)$ \code{else} \code{not} $p(a, b, c)$ adalah \textbf{kalimat}, karena:
  \begin{itemize}
    \item $(\forall x)\, p(a, b, x)$ dan $(\exists y)\, q(y)$ adalah kalimat
    \item \code{not} $p(a, b, c)$ adalah kalimat
    \item Maka kondisional \code{if} kalimat \code{then} kalimat \code{else} kalimat adalah kalimat
  \end{itemize}
\end{enumerate}
\end{contoh}

\subsection{Definisi Ekspresi}

\begin{konsep}
Suatu \logic{ekspresi} dalam kalkulus predikat dapat berupa kalimat atau term.
\end{konsep}

\begin{contoh}
\begin{center}
\begin{tabular}{ll}
  $x$ & merupakan ekspresi (term) \\
  $f(x, y)$ & merupakan ekspresi (term) \\
  $(\exists x)\, p(x)$ & merupakan ekspresi (kalimat) \\
\end{tabular}
\end{center}
\end{contoh}

\subsection{Subterm, Subkalimat, dan Subekspresi}

\begin{konsep}
\begin{enumerate}
  \item \logic{Subterm} dari term $t$ atau dari kalimat $A$ adalah setiap term antara yang digunakan untuk membangun $t$ atau $A$
  \item \logic{Subkalimat} adalah setiap kalimat antara yang digunakan untuk membangun kalimat yang lebih luas
  \item \logic{Subekspresi} adalah subterm atau subkalimat yang terdapat pada sebuah ekspresi
\end{enumerate}
\end{konsep}

\begin{contoh}
Sebutkan semua subterm dan subkalimat dari ekspresi:
$$E : \text{\code{if}} \; (\forall x)\, q(x, f(a)) \; \text{\code{then}} \; f(a) \; \text{\code{else}} \; b$$

\begin{itemize}
  \item \textbf{Subterm}: $a$, $x$, $f(a)$, $b$, dan $E$ sendiri (karena $E$ adalah term)
  \item \textbf{Subkalimat}: $q(x, f(a))$ dan $(\forall x)\, q(x, f(a))$
  \item Semuanya merupakan \textbf{subekspresi} dari $E$
\end{itemize}
\end{contoh}
